% ============================================================
\section{Introduction}
\label{sec:intro}
% ============================================================

Quantum machine learning has generated substantial interest around the
hypothesis that quantum-inspired data representations could improve learning
on classical hardware~\citep{havlicek2019,schuld2019}.
Amplitude, angle, and basis encoding are the canonical strategies for loading
classical data into quantum circuits, and they induce specific geometric
transformations: $\ell_2$-normalisation with power-of-two dimensional
expansion, per-feature trigonometric lifting, and binary discretisation
respectively.
Applied as explicit classical feature maps, these transformations can be
evaluated independently of quantum hardware, making it possible to isolate the
geometric contribution of the encoding from hardware-specific quantum effects.

Prior work has shown that classical algorithms can match or exceed quantum
algorithms on certain linear-algebra tasks~\citep{tang2019} and that classical
kernels can be trained to match quantum kernel performance on benchmark
datasets~\citep{kubler2021,huang2021}.
However, these comparisons focus on accuracy endpoints.
The question of \emph{why} QIE representations succeed or fail, in terms of
rank structure, spectral conditioning, and kernel alignment, has not been
systematically addressed.

This paper makes the following contributions:
\begin{enumerate}
  \item A ten-dataset, five-seed benchmark of amplitude, angle, and basis
        encoding against eight classical baselines (raw linear, RBF-SVM,
        polynomial degree-2 and degree-3, Random Fourier Features, PCA,
        and two MLP variants) under matched representational budget.
  \item A spectral attribution framework that identifies the precise failure
        mode of each encoding: rank collapse, geometric equivalence, and
        Hamming geometry mismatch.
  \item A falsifiable prediction: QIE approaches competitive performance when
        the dataset's intrinsic geometry uniformly fills the encoding's output
        space.  We verify this prediction on the high-rank noise dataset.
  \item An overhead characterisation showing that encoding cost is not the
        bottleneck (amplitude: $35\,\mathrm{ms}$, matching raw standardisation;
        polynomial degree-2: $10.98\,\mathrm{s}$).
\end{enumerate}

The paper is organised as follows.
Section~\ref{sec:related} surveys related work.
Section~\ref{sec:methods} describes the encodings, baselines, datasets,
and evaluation protocol.
Section~\ref{sec:results} reports predictive performance and statistical tests.
Section~\ref{sec:analysis} presents the spectral attribution analysis.
Section~\ref{sec:overhead} characterises practical overhead.
Section~\ref{sec:discussion} discusses implications and limitations.
Section~\ref{sec:conclusion} concludes.
